\chapter{p-adic Numbers}\label{p-adic-numbers}
\chapterstatus{planned}

Absolute values on $\QQ$, the completions $\QQ_p$, and the local-global principle for quadratic forms.

\noindent\textbf{Prerequisites.} Chapter~\ref{number-theory} (Elementary Number Theory), Chapter~\ref{real-analysis} (Real Analysis), Chapter~\ref{general-topology} (General Topology), Chapter~\ref{geometry-of-numbers} (Geometry of Numbers) and Chapter~\ref{analytic-number-theory} (Analytic Number Theory).

\noindent\textbf{Planned contents.}
\begin{enumerate}
\item Absolute values; Ostrowski's theorem; the product formula.
\item Construction of $\QQ_p$ and $\ZZ_p$ (completion and inverse limit), $p$-adic expansions, the topology of $\ZZ_p$.
\item Hensel's lemma and its applications: squares in $\QQ_p$, roots of unity, the structure of $\QQ_p^\times$.
\item $p$-adic analysis: series, $\exp$ and $\log$, Mahler's theorem, Strassmann's theorem; the $p$-adic proof of the Skolem--Mahler--Lech theorem (statement).
\item The Hilbert symbol and Hilbert reciprocity; quadratic forms over $\QQ_p$; the Hasse--Minkowski theorem (its proof uses Dirichlet's theorem on primes in progressions); Legendre's theorem on $ax^2 + by^2 + cz^2 = 0$ again.
\item Exercises with solutions.
\end{enumerate}
